\section{Conclusion}
  \label{sec:conclusion}

  This paper analysed the dynamical mechanism governing the projective resolution
  along the relaxation cascade.
  The companion paper~\cite{SpectralStratigraphy} established that the
  representation-theoretic structure of ADE binary Cayley graphs fixes the
  number and order of stratigraphic levels, while the dynamical growth of
  $\Lambda_{\mathrm{proj}}(n)$ controls their separation.
  The present paper addresses this second question.

  The main positive results are: the isoperimetric law
  $\Lambda_{\mathrm{proj}}(n) \asymp h(G(n))^2$
  (Proposition~\ref{prop:Lambda-h}); the Cheeger enclosure
  $\lambda_2^2 \lesssim \Lambda_{\mathrm{proj}} \lesssim \lambda_2$
  (Corollary~\ref{cor:cheeger-enclosure}); the linear law
  $\Lambda_{\mathrm{proj}} \asymp \lambda_2$ under spectral-isoperimetric
  saturation (Corollary~\ref{cor:linear-law}); the linear growth
  $N(\lambda; n) \sim F_{\mathrm{KM}}(\lambda) \cdot n$
  (Proposition~\ref{prop:Nproj-linear}); and the factorised mass ratio formula
  \[
    \frac{\mathcal{M}_i}{\mathcal{M}_j}
    \;=\;
    \frac{F_{\mathrm{KM}}(\lambda_i)}{F_{\mathrm{KM}}(\lambda_j)}
    \cdot
    \frac{\dim\rho_{\lambda_j}}{\dim\rho_{\lambda_i}},
  \]
  in which the structural identity $F_{\mathrm{KM}}(1) = 1/2$ holds universally.

  Two limitations are established: the saturation mechanism inverts the mass
  ordering expected from the admissibility envelope, and the predicted mass
  ratios lie in $[0.10,\, 2.2]$, far below the observed $10^{-3}$--$10^{-5}$
  range.
  Both limitations are traced to the asymptotic constancy of
  $\Lambda_{\mathrm{proj}}$ in the LPS model at fixed $p$.

  The primary open direction is direction~(O1): allowing $p(n) \to \infty$ with
  $n$, so that the Kesten--McKay support narrows progressively, the ADE levels
  leave the support in a rank-dependent order, and the admissibility and
  saturation conditions are coupled.
  This regime, in which the effective relational valence of the substrate
  increases along the cascade, is the natural candidate for recovering both the
  correct mass ordering and the observed amplification, and will be the subject
  of the next paper in the admissibility sub-programme.
